\documentclass[10pt]{article}
\usepackage[T1]{fontenc}
\usepackage[utf8]{inputenc}
\usepackage[french]{babel}
\usepackage{amsmath,amssymb,mathrsfs,mathpazo}
\usepackage{geometry}
\geometry{a4paper,margin=2cm}
\usepackage{graphicx,array,tabularx,tabularray,colortbl,multirow,multicol,booktabs,enumitem,pifont,ulem,fancybox,diagbox,cancel}
\usepackage{pstricks,pstricks-add,pst-plot,pst-tree,pst-node,pst-text,pst-func,pst-math,pst-bspline,pst-3dplot,pst-eucl,pst-solides3d}
\let\pscalculate\undefined
\usepackage{pst-bezier}
\usepackage{tikz,pgfplots,tkz-tab}
\usetikzlibrary{arrows,arrows.meta,patterns,calc,positioning,trees,decorations.markings,shapes.geometric,shapes.misc,shapes.arrows,intersections,decorations.pathmorphing,decorations.pathreplacing,backgrounds,shadows,babel}
\pgfplotsset{compat=1.18}
\usepgfplotslibrary{fillbetween}
\usepackage[np]{numprint}
\usepackage{listings,makecell,pas-tableur,twoopt}
\let\barre\undefined
\usepackage[pstricks]{tablvar}
\definecolor{darkgreen}{RGB}{0,100,0}
\providecommand{\red}{\color{red}}
\providecommand{\blue}{\color{blue}}
\providecommand{\green}{\color{green}}
\providecommand{\black}{\color{black}}
\providecommand{\white}{\color{white}}
\providecommand{\gray}{\color{gray}}
\providecommand{\cyan}{\color{cyan}}
\providecommand{\magenta}{\color{magenta}}
\providecommand{\yellow}{\color{yellow}}
\providecommand{\violet}{\color{violet}}
\providecommand{\orange}{\color{orange}}
\providecommand{\euro}{\texteuro}
\providecommand{\EUR}{\texteuro}
\providecommand{\ieme}{\textsuperscript{e}}
\providecommand{\degres}{\ensuremath{^\circ}}
\providecommand{\degree}{\ensuremath{^\circ}}
\providecommand{\vv}[1]{\overrightarrow{#1}}
\providecommand{\Coord}[1]{#1}
\providecommand{\decosix}{}
\providecommand{\starredbullet}{\ensuremath{\bullet}}
\pagestyle{empty}
\setlength{\parindent}{0pt}
\usepackage{adjustbox}\DeclareUnicodeCharacter{2236}{:}\DeclareUnicodeCharacter{2190}{\ensuremath{\leftarrow}}
\providecommand{\R}{\mathbb{R}}
\providecommand{\N}{\mathbb{N}}
\providecommand{\Z}{\mathbb{Z}}
\providecommand{\Q}{\mathbb{Q}}
\providecommand{\D}{\mathbb{D}}
\providecommand{\C}{\mathbb{C}}
\providecommand{\e}{{\mathrm{e}}}
\providecommand{\vect}[1]{\overrightarrow{#1}}
\providecommand{\vectt}[1]{\overrightarrow{\mathrm{#1}}}
\providecommand{\vvt}[1]{\overrightarrow{\mathrm{#1}}}
\providecommand{\np}[2][]{{#2~#1}}
\providecommand{\vv}[1]{{\overrightarrow{#1}}}
\renewcommand{\d}{{\mathrm{d}}}
\providecommand{\degres}{{^\circ}}
\providecommand{\degre}{{^\circ}}
\providecommand{\blue}{\color{blue}}
\providecommand{\red}{\color{red}}
\providecommand{\enskip}{\quad}
\providecommand{\llbracket}{\lbrack\!\lbrack}
\providecommand{\rrbracket}{\rbrack\!\rbrack}
\providecommand{\euro}{{\text{€}}}
\providecommand{\no}{{\text{n°}}}
\providecommand{\interval}[2]{{[#1;#2]}}
\providecommand{\intervFO}[2]{{[#1;#2[}}
\providecommand{\dotfill}{\ldots}

\providecommand{\ds}{\displaystyle}
\providecommand{\textding}[1]{#1}
\providecommand{\Oij}{\(\mathrm{O};\vec{\imath},\vec{\jmath}\)}
\providecommand{\Oijk}{\(\mathrm{O};\vec{\imath},\vec{\jmath},\vec{k}\)}
\providecommand{\Ouv}{\(\mathrm{O};\vec{u},\vec{v}\)}
\providecommand{\intervFF}[2]{\ensuremath{[#1;#2]}}
\renewcommand{\euro}{\ifmmode\text{\texteuro}\else\texteuro\fi}
\DeclareUnicodeCharacter{00B0}{\ifmmode^{\circ}\else\textdegree\fi}
\DeclareUnicodeCharacter{20AC}{\ifmmode\text{\texteuro}\else\texteuro\fi}
% Énoncés : Éducation nationale. Transcription et corrigés : APMEP, Denis Vergès et contributeurs. Découpage et classement : CoopMaths, Sébastien Lozano, Éric Elter et contributeurs. Adaptation de la présentation : Maths Explorer.
% CC BY-SA : https://coopmaths.fr/www/about/
% https://coopmaths.fr/alea/static/e3c/2020/tex/e3c_2020_00_sujet24_4_cor.tex
\begin{document}
\noindent\textit{APMEP / CoopMaths. CC BY-SA. Présentation adaptée par Maths Explorer.}\par\medskip
\subsection*{Question 1}

On a :
\begin{align*}
	f_1 &= f_0 + 108 \\
	&= 2500 + 108 = 2608.
\end{align*}

Augmenter de 3,8\%, c'est multiplier par :
\begin{align*}
	1 + \frac{3,8}{100} &= 1 + 0,038 = 1,038.
\end{align*}
Ainsi, 
\begin{align*}
	c_1 &= c_0 \times 1,038 = 2500 \times 1,038 = 2595.
\end{align*}

\subsection*{Question 2}

La suite $(f_n)$ est une suite arithmétique de raison $108$ et de premier terme $f_0 = 2500$, donc :
\begin{align*}
	f_{n+1} &= f_n + 108.
\end{align*}

La suite $(c_n)$ est une suite géométrique de raison $1,038$ et de premier terme $c_0 = 2500$, donc :
\begin{align*}
	c_{n+1} &= c_n \times 1,038.
\end{align*}

\subsection*{Question 3}

On a les expressions suivantes pour $n \in \mathbb{N}$ :
\begin{align*}
	f_n &= 2500 + 108n, \\
	c_n &= 2500 \times 1,038^n.
\end{align*}

\subsection*{Question 4}

On complète l'algorithme pour déterminer le nombre de mois à attendre pour que le nombre potentiel de flacons commandés dépasse celui des flacons produits :
\begin{lstlisting}[style=python]
	n = 0
	f = 2500
	c = 2500
	
	while f > c:
		n = n + 1
		f = f + 108
		c = c * 1.038
\end{lstlisting}

\subsection*{Question 5}

Il faut comparer $F_{12}$ et $C_{12}$, les sommes cumulées des flacons produits et commandés après 12 mois.

On calcule $F_{12}$ :
\begin{align*}
	F_{12} &= f_0 + f_1 + \dots + f_{11} \\
	&= 12 \times \left( \frac{f_0 + f_{11}}{2} \right) \\
	&= 12 \times \left( \frac{2500 + (2500 + 11 \times 108)}{2} \right) \\
	&= 12 \times (2500 + 594) = 12 \times 3094 = 37128.
\end{align*}

Pour $C_{12}$, on a une suite géométrique de raison $1,038$ :
\begin{align*}
	C_{12} &= c_0 + c_1 + \dots + c_{11}.
\end{align*}
En utilisant la somme des termes d'une suite géométrique, on a :
\begin{align*}
	C_{12} &= 2500 \times \frac{1,038^{12} - 1}{1,038 - 1} \\
	&\approx 37136.
\end{align*}

Ainsi, à la fin de l'année, le nombre total de flacons commandés ($37136$) sera légèrement supérieur à celui des flacons produits ($37128$), ce qui confirme que le modèle prédit une insuffisance de production.

\end{document}
